% GENERATED FILE. Edit canonical lessons, then run npm run generate:books.
% canonical-source-hash: fnv1a64:c019747e7a9608bd
% canonical-lessons: KA-C50-now, KA-C50-day-after, KA-C50-journey, KA-C50-set-out, KA-C50-farewell

\chapter{Taking Your Leave}
\label{ch:ka-leave-more}
\hlchaptermodality{50}

\begin{chapteropening}
\textbf{By the end of this chapter:} \emph{I can take my leave in Kannada and say when I will be back, and explain why the word for the farewell itself is built on the verb for giving.}

Complete KA-C50-farewell and say all five words in order.
\end{chapteropening}

\section[īga]{\kn{ಈಗ} (īga) — now}

\label{lesson:KA-C50-now}



\begin{quote}
Before the new run: what did \emph{hāge} mean, and what did \emph{bahuśaḥ} mean?
\end{quote}

\subsection*{You'll want to know: \kn{ಈಗ}}
\textbf{\kn{ಈಗ}} (\emph{īga}) — \textquotedblleft{}now\textquotedblright{}.

\kn{ಈಗ} is 'now', and it stands on the same \kn{ಇ}- of nearness the chapter before ended on. \kn{ಆಗ} (\emph{āga}) is its partner: 'then, at that time', carrying the \kn{ಅ}- of distance.

So Kannada names this moment and that moment by the same means it names this thing and that thing. Time is handled as one more place to point at.

\begin{grammarlens}[title={What you've built}]
The first of five words for going.
\end{grammarlens}

\subsection*{Your turn}
Say these aloud:

\begin{itemize}
\raggedright
  \item \emph{īga}
  \item it once more, slowly
  \item \emph{īga}, then \emph{āga}, and point at the near one and the far one
\end{itemize}

\begin{tcolorbox}[breakable,colback=teal!4,colframe=teal!35!black,title={Before you move on}]
What does \kn{ಈಗ} mean? (\textquotedblleft{}now\textquotedblright{}.) Say it once more.
\end{tcolorbox}
\section[nāḍiddu]{\kn{ನಾಡಿದ್ದು} (nāḍiddu) — the day after tomorrow}

\label{lesson:KA-C50-day-after}



\begin{quote}
Before the new one: what did \emph{īga} mean?
\end{quote}

\subsection*{You'll want to know: \kn{ನಾಡಿದ್ದು}}
\textbf{\kn{ನಾಡಿದ್ದು}} (\emph{nāḍiddu}) — \textquotedblleft{}the day after tomorrow\textquotedblright{}.

English needs four words for this. Kannada needs one, and keeps a matching one running the other way: \kn{ಮೊನ್ನೆ} (\emph{monne}), the day before yesterday.

You can hear \kn{ನಾಳೆ}, 'tomorrow', at the front of \kn{ನಾಡಿದ್ದು} — the two belong together, one step out and two steps out. A language that sows and marries by the calendar finds it worth keeping single words this far from today.

\begin{grammarlens}[title={What you've built}]
Two: now, and the day after tomorrow.
\end{grammarlens}

\subsection*{Your turn}
Say these aloud:

\begin{itemize}
\raggedright
  \item \emph{nāḍiddu}
  \item it once more, slowly
  \item \emph{īga}, \emph{nāḷe}, \emph{nāḍiddu} — today's moment, then one day out, then two
\end{itemize}

\begin{tcolorbox}[breakable,colback=teal!4,colframe=teal!35!black,title={Before you move on}]
What does \kn{ನಾಡಿದ್ದು} mean? (\textquotedblleft{}the day after tomorrow\textquotedblright{}.) Say it once more.
\end{tcolorbox}
\section[prayāṇa]{\kn{ಪ್ರಯಾಣ} (prayāṇa) — a journey}

\label{lesson:KA-C50-journey}



\begin{quote}
Before the new one: what did \emph{nāḍiddu} mean?
\end{quote}

\subsection*{You'll want to know: \kn{ಪ್ರಯಾಣ}}
\textbf{\kn{ಪ್ರಯಾಣ}} (\emph{prayāṇa}) — \textquotedblleft{}a journey\textquotedblright{}.

Sanskrit \emph{pra-yāṇa}: \emph{pra}, 'forth, forward', laid on a root meaning to go. A going-forth.

That \emph{pra-} is very old and very widely travelled. It is the \emph{pro-} of Latin, the \emph{pro-} of Greek and the \emph{forth} of English, one prefix that four languages inherited rather than borrowed from one another. Wish someone \kn{ಶುಭ} \kn{ಪ್ರಯಾಣ}, a good journey, and half the phrase is a word this book taught you on a quite different errand.

\begin{grammarlens}[title={What you've built}]
Three: now, the day after tomorrow, a journey.
\end{grammarlens}

\subsection*{Your turn}
Say these aloud:

\begin{itemize}
\raggedright
  \item \emph{prayāṇa}
  \item it once more, slowly
  \item \emph{śubha prayāṇa} to someone leaving, then \emph{prayāṇa} on its own
\end{itemize}

\begin{tcolorbox}[breakable,colback=teal!4,colframe=teal!35!black,title={Before you move on}]
What does \kn{ಪ್ರಯಾಣ} mean? (\textquotedblleft{}a journey\textquotedblright{}.) Say it once more.
\end{tcolorbox}
\section[horaḍu]{\kn{ಹೊರಡು} (horaḍu) — to set out}

\label{lesson:KA-C50-set-out}



\begin{quote}
Before the new one: what did \emph{prayāṇa} mean?
\end{quote}

\subsection*{You'll want to know: \kn{ಹೊರಡು}}
\textbf{\kn{ಹೊರಡು}} (\emph{horaḍu}) — \textquotedblleft{}to set out\textquotedblright{}.

\kn{ಹೊರಡು} is to set out, to start off, to leave a place with somewhere else in mind — not merely to go, for which Kannada already handed you a word.

Its front half is \kn{ಹೊರ} (\emph{hora}), 'outside'. To set out is to go outward. And \kn{ಹೊರ} is one more inherited word with an \emph{h} where Tamil has a \emph{p}: \ta{புறம்} (\emph{puṟam}) is the outside there. The law turns up again in the middle of an ordinary travelling verb.

\begin{grammarlens}[title={What you've built}]
Four. Now, the day after tomorrow, a journey, and setting out on it.
\end{grammarlens}

\subsection*{Your turn}
Say these aloud:

\begin{itemize}
\raggedright
  \item \emph{horaḍu}
  \item it once more, slowly
  \item \emph{horaḍu}, then \emph{hōgu}, and say which of the two has a destination in it
\end{itemize}

\begin{tcolorbox}[breakable,colback=teal!4,colframe=teal!35!black,title={Before you move on}]
What does \kn{ಹೊರಡು} mean? (\textquotedblleft{}to set out\textquotedblright{}.) Say it once more.
\end{tcolorbox}
\section[bīḷkoḍuge]{\kn{ಬೀಳ್ಕೊಡುಗೆ} (bīḷkoḍuge) — a send-off, a farewell}

\label{lesson:KA-C50-farewell}



\begin{quote}
Before the new one: what did \emph{horaḍu} mean?
\end{quote}

\subsection*{You'll want to know: \kn{ಬೀಳ್ಕೊಡುಗೆ}}
\textbf{\kn{ಬೀಳ್ಕೊಡುಗೆ}} (\emph{bīḷkoḍuge}) — \textquotedblleft{}a send-off, a farewell\textquotedblright{}.

Take it apart and something you already own falls out of it. \kn{ಬೀಳ್ಕೊಡು} is 'to see someone off', and the \kn{ಕೊಡು} sitting in the middle of it is the plain verb 'give' from the chapter of commands. The ending \emph{-ge} turns the whole action into a thing: a send-off.

Notice whose act it is. Kannada makes the leave-taking belong to the one who stays: the host gives the departure away. The traveller only walks.

Five words for going: now, the day after tomorrow, a journey, to set out, and the send-off at the end of it.

\begin{grammarlens}[title={What you've built}]
Five words for going, and a way to end a visit without repeating yourself.
\end{grammarlens}

\subsection*{Your turn}
Say these aloud:

\begin{itemize}
\raggedright
  \item \emph{bīḷkoḍuge}
  \item it once more, slowly
  \item all five in order, then \emph{bīḷkoḍuge} and the \emph{koḍu} hiding inside it
\end{itemize}

\begin{tcolorbox}[breakable,colback=teal!4,colframe=teal!35!black,title={Before you move on}]
What does \kn{ಬೀಳ್ಕೊಡುಗೆ} mean? (\textquotedblleft{}a send-off, a farewell\textquotedblright{}.) Say it once more.
\end{tcolorbox}
